%%%%%%%% Runs kept out of the main report (paper.tex). To become the ICML paper's appendix.
%%%%%%%% Terminology follows writeup/glossary.md. All results: validation split, one seed (s42) unless noted.
\documentclass[11pt]{article}
\usepackage[T1]{fontenc}
\usepackage[margin=2.2cm]{geometry}
\usepackage{booktabs}
\usepackage{graphicx}
\usepackage{amsmath}
\usepackage{amssymb}
\usepackage{hyperref}
\usepackage[capitalize,noabbrev]{cleveref}
\newcommand{\bok}{best-of-$K$}
\newcommand{\Bok}{Best-of-$K$}
\newcommand{\Y}{|Y(x)|}

\title{Extra Results\\\large Runs outside the main report}
\date{Draft, 23 November 2026}

\begin{document}
\maketitle

Everything here uses the setup of the main report: GRAM's network, 150k optimiser steps on 768 sequences per step,
100 samples per input, validation inputs, run summaries (mean of the last four scored checkpoints) unless a table says
final checkpoint. Values are validity (\%) / unique-solution validity (\%) / distinct@100 unless stated.

\section{N-Queens}

All 724 solutions of 10-Queens with 7, 8 or 9 queens removed; each distinct partial board is one input. 43{,}420
inputs, split 36{,}907 training / 1{,}000 validation (the first 1000 test inputs) / 5{,}513 test. Solutions per input:
1--93. N-Queens inputs have few solutions, and 100 GRAM samples reach all of them, so coverage@10 is reported too.

\begin{table}[h]
\centering\small
\resizebox{\textwidth}{!}{\begin{tabular}{@{}lccccc@{}}
\toprule
Model & Validity & Unique-sol.\ validity & Distinct@100 & Coverage@10 & Posterior sampling validity \\
\midrule
GRAM $\beta=0.0225$ & 14.2 & 7.7 & 3.4 & 0.42 & 100.0 \\
GRAM $\beta=0.045$ (reference) & 88.9 & 79.9 & 3.7 & 0.89 & 99.7 \\
GRAM $\beta=0.09$ & 2.3 & 2.8 & 0.2 & 0.04 & 1.7 \\
\Bok{} ($K=8$) & 97.1 & 94.2 & 3.0 & 0.68 & -- \\
\quad state noise $\sigma=0.5$ / 1 / 2 & 97.9 / 97.3 / 93.8 & & 3.2 / 3.5 / 3.6 & & \\
\bottomrule
\end{tabular}}
\caption{N-Queens. Coverage@10 and posterior sampling at the final checkpoint; state noise rows give validity /
distinct@100 on the final checkpoint.}
\end{table}

At half the reference $\beta$, validity falls as training continues (47\% at 15k steps, 14\% at 150k) while posterior
sampling stays 100\% valid: prior--posterior mismatch. At twice the reference $\beta$, unique-solution validity also
falls (2.8\%) and posterior sampling fails too (1.7\%), which posterior collapse does not explain: the network itself
stops producing valid boards. The baseline's coverage stops rising long before every solution (0.81 at 50 samples,
0.84 at 100), while GRAM reaches 1.00 at 100. Unseen solutions: on the 14 training inputs with more than 16
solutions, the unseen share is 35.5\% for GRAM and 38.7\% for \bok, against 35.3\% and 35.6\% for a uniform sampler.

\section{GRAM: further $\beta$ values and the maze by number of solutions}

\begin{table}[h]
\centering\small
\resizebox{\textwidth}{!}{\begin{tabular}{@{}llcl@{}}
\toprule
Task & $\beta$ & Validity / unique / distinct@100 & Training KL (end) \\
\midrule
Maze & 0.005 ($\frac14$) & 9.5 / 9.3 / 5.3 & \\
 & 0.02 (reference) & 35.9 / 97.8 / 0.8 & 0.1--0.5 \\
 & 0.04 (2$\times$) & 33.8 / 98.6 / 0.0 & $\sim$0.001 \\
 & 0.08 (4$\times$) & 33.6 / 99.3 / 0.0 & $\sim$0.0003 \\
Graph col. & 0.225 & 71.5 / 93.4 / 31.2 & 3.4 \\
 & 0.9 & 86.5 / 100.0 / 2.2 & 0.0004 \\
Sudoku & 0.1 (GRAM's Sudoku-Extreme value, $2\times$ ours) & 35.4 / 98.3 / 0.3 & \\
\bottomrule
\end{tabular}}
\caption{Extra $\beta$ points and training KL (before multiplying by $\beta$), read off training curves; blank where
not read. Maze posterior sampling at $\frac14$: 94\% valid.}
\end{table}

\begin{table}[h]
\centering\small
\begin{tabular}{@{}lcccccc@{}}
\toprule
Maze solutions & 1 & 2--3 & 4--15 & 16--63 & 64--255 & 256+ \\
\midrule
Inputs & 660 & 144 & 300 & 299 & 297 & 300 \\
GRAM, $\beta=0.02$ & 97.8 & 48.4 & 19.2 & 1.9 & 0.1 & 0.0 \\
\Bok & 75.9 & 79.2 & 90.9 & 97.0 & 97.0 & 95.2 \\
\bottomrule
\end{tabular}
\caption{Maze validity (\%) by number of solutions, first 2000 training inputs, final checkpoints.}
\end{table}

\paragraph{$\beta$ from data statistics (Q1).} A per-task $\beta$ computed from the stored solutions alone (the
per-cell entropy over stored solutions per labelled cell, divided by $\ln$ of the number stored) does not predict the
best $\beta$: the ratio of best $\beta$ to its median is 0.21 (maze), 0.52 (N-Queens), 0.44 (Sudoku) and 1.69 (graph
colouring), an 8-fold spread, worse than the rule $\beta\times\text{cells}\approx4.5$.

\paragraph{Seeds.} Second seeds (s43): GRAM Sudoku 0.05: 93.1 / 98.0 / 34.4 (s42: 92.5 / 97.6 / 34.0); \bok{} Sudoku:
99.2 / -- / 9.4 (s42: 99.2 / 99.3 / 10.4); GRAM maze 0.01: 62.6 / 44.6 / 36.2 (s42: 67.3 / 48.2 / 35.3). Seed spread is
under 1 point on Sudoku and about 5 points of validity on mazes. Both GRAM seeds dip mid-training and recover.

\paragraph{State noise on GRAM.} Sudoku, $\beta=0.05$, $\sigma=0.5$: 91.2 / 97.2 / 33.7, against 92.5 / 97.6 / 34.0
without it.

\section{\Bok{}: number of rollouts and larger state noise}

\begin{table}[h]
\centering\small
\resizebox{\textwidth}{!}{\begin{tabular}{@{}llcccc@{}}
\toprule
Task & & $K=1$ (MinLoss) & $K=4$ & $K=8$ & $K=16$ \\
\midrule
& Inputs per optimiser step & 768 & 192 & 96 & 48 \\
\midrule
Sudoku & Validity / unique / distinct@100 & 99.5 / 99.8 / 2.5 & 99.8 / -- / 5.3 & 99.2 / 99.3 / 10.4 & 99.1 / -- / 7.8 \\
 & Distinct@100, $\sigma=0.5$ & 18.0 & & 27.7 & \\
Maze & Validity / unique / distinct@100 & 88.4 / 94.1 / 1.0 & & 83.3 / 75.4 / 2.2 & 13.4 / 20.5 / 1.9 \\
\bottomrule
\end{tabular}}
\caption{The baseline across $K$ at matched compute. Sudoku $K=16$ trained for 300k steps: 98.6 / -- / 12.9. Maze
$K=16$ second seed: 40.7 / 43.0 / 1.6.}
\end{table}

On Sudoku, \bok{} adds diversity that MinLoss does not have, and state noise adds about the same on top of either
($+15.5$ at $K=1$, $+17.3$ at $K=8$). On mazes, larger $K$ lowers unique-solution validity, and $K=16$ fails in every
group of solution counts (5--21\%), with checkpoints swinging between 9 and 78\% on both seeds. Since $K$ also sets the
inputs per optimiser step, these runs do not separate the number of rollouts from the number of inputs behind each step.

\begin{table}[h]
\centering\small
\begin{tabular}{@{}lcccccc@{}}
\toprule
State noise $\sigma$ & 0 & 0.25 & 0.5 & 1 & 1.5 & 2 \\
\midrule
Sudoku & 99.2 / 10.4 & 99.5 / 20.9 & 99.5 / 27.7 & 97.5 / 29.1 & 23.0 / 1.7 & 0.0 / 0.0 \\
Graph col. & 99.8 / 38.4 & & 99.6 / 39.5 & 99.4 / 40.4 & & 66.8 / 29.5 \\
Maze & 87.7 / 2.1 & & 88.5 / 3.2 & 89.9 / 4.4 & & \\
\bottomrule
\end{tabular}
\caption{\Bok{} ($K=8$) with state noise: validity / distinct@100, final checkpoint. Validity breaks between
$\sigma=1$ and 1.5 on Sudoku.}
\end{table}

\section{Other ways to add randomness to the baseline}

Each of these was tried to raise the baseline's diversity without a posterior. None beats \bok{} with state noise.
On Sudoku, every source of randomness tried tops out at about 28 distinct@100 at 99\% validity (GRAM: 34.0 at 92.5\%).

\begin{table}[h]
\centering\small
\resizebox{\textwidth}{!}{\begin{tabular}{@{}llccc@{}}
\toprule
Variant & What changes & Sudoku & Graph col. & Maze \\
\midrule
Sample-level stochastic guidance & one $\mathcal N(0,I)$ vector per sample, not per position & 99.7 / 99.7 / 8.9 & 47.7 / 58.9 / 15.0 & 85.3 / 87.2 / 7.0 \\
STRM ($K=8$) & random initial high-level state, no stochastic guidance & 99.9 / 99.9 / 16.8 & 99.9 / -- / 19.9 & \\
\quad $\sigma=0.5$ / 1 / 1.5 & state noise on STRM (final checkpoint) & 99.9 / 20.1;\ 97.9 / 26.1;\ 21.8 / 0.9 & & \\
STRM + stochastic guidance & both & 94.8 / -- / 19.2 & & \\
STRM ($K=1$) & implementation check & 99.5 / 100.0 / 1.0 & & \\
Training state noise $\sigma=1$ & state noise also in training (Sudoku at 58k steps) & see text & see text & \\
\bottomrule
\end{tabular}}
\caption{Validity / unique-solution validity / distinct@100 (state-noise rows: validity / distinct@100). Blank: not
run. STRM on graph colouring was unstable across checkpoints.}
\end{table}

\paragraph{Training state noise.} Training \bok{} with state noise $\sigma=1$ moves the point at which test-time state
noise breaks validity, but not the trade-off: on Sudoku at 58k steps, test $\sigma=0$ / 1 / 1.5 / 2 give 99.3 / 4.1,
99.5 / 21.4, 99.0 / 27.8 and 94.0 / 27.2, the same best point as plain \bok{} at $\sigma=0.5$ (99.5 / 27.7). On graph
colouring at 80k steps validity broke (36.7 / 20.0 at $\sigma=1$). Both runs were stopped.

\paragraph{STRM on heyawake.} 86.8 / 83.2 / 4.1, run summary. State noise on it: pending.

\section{Unseen solutions at the tuned maze $\beta$}

GRAM at $\beta=0.01$ on 881 maze training inputs with more than 16 solutions: unseen share 79.2\% (uniform 78.9\%),
43.1 distinct unseen per input, against 0.2 at the reference $\beta$.

\section{Ratio to uniform}

$d$ divided by $\Y(1-(1-1/\Y)^v)$, the expected number of distinct answers in $v$ uniform draws. It conditions on the
number of valid samples, so it approaches 1 for any model with few valid samples, and we never use it to compare
methods. At the reference $\beta$: maze 0.84 (GRAM) and 0.14 (\bok), N-Queens 1.00 and 0.86, Sudoku 0.97 and 0.44,
graph colouring 0.96 and 0.81, heyawake 0.95 (GRAM). On mazes GRAM's 0.84 reflects its few valid samples, not
diversity: its distinct@100 is 0.8 against the baseline's 2.2.

\section{Random 3-SAT}

A benchmark of random 3-SAT formulas over 16 variables (clause density $U[1,4.5]$, up to a quarter of the variables
given) was built the same way. The baseline did not learn it: 0.4\% of its samples were valid after 50k optimiser
steps, close to random assignments. GRAM failed the same way in an earlier configuration. Excluded until the cause is
understood.

\section{Light up}

A light up (akari) benchmark from Pencil Puzzle Bench was built (builder and checker in the code) but not used: five
tasks were fixed before it finished.

\end{document}
